\documentclass[10pt]{article}
\usepackage[T1]{fontenc}
\usepackage[utf8]{inputenc}
\usepackage[french]{babel}
\usepackage{amsmath,amssymb,mathrsfs,mathpazo}
\usepackage{geometry}
\geometry{a4paper,margin=2cm}
\usepackage{graphicx,array,tabularx,tabularray,colortbl,multirow,multicol,booktabs,enumitem,pifont,ulem,fancybox,diagbox,cancel}
\usepackage{pstricks,pstricks-add,pst-plot,pst-tree,pst-node,pst-text,pst-func,pst-math,pst-bspline,pst-3dplot,pst-eucl,pst-solides3d}
\let\pscalculate\undefined
\usepackage{pst-bezier}
\usepackage{tikz,pgfplots,tkz-tab}
\usetikzlibrary{arrows,arrows.meta,patterns,calc,positioning,trees,decorations.markings,shapes.geometric,shapes.misc,shapes.arrows,intersections,decorations.pathmorphing,decorations.pathreplacing,backgrounds,shadows,babel}
\pgfplotsset{compat=1.18}
\usepgfplotslibrary{fillbetween}
\usepackage[np]{numprint}
\usepackage{listings,makecell,pas-tableur,twoopt}
\let\barre\undefined
\usepackage[pstricks]{tablvar}
\definecolor{darkgreen}{RGB}{0,100,0}
\providecommand{\red}{\color{red}}
\providecommand{\blue}{\color{blue}}
\providecommand{\green}{\color{green}}
\providecommand{\black}{\color{black}}
\providecommand{\white}{\color{white}}
\providecommand{\gray}{\color{gray}}
\providecommand{\cyan}{\color{cyan}}
\providecommand{\magenta}{\color{magenta}}
\providecommand{\yellow}{\color{yellow}}
\providecommand{\violet}{\color{violet}}
\providecommand{\orange}{\color{orange}}
\providecommand{\euro}{\texteuro}
\providecommand{\EUR}{\texteuro}
\providecommand{\ieme}{\textsuperscript{e}}
\providecommand{\degres}{\ensuremath{^\circ}}
\providecommand{\degree}{\ensuremath{^\circ}}
\providecommand{\vv}[1]{\overrightarrow{#1}}
\providecommand{\Coord}[1]{#1}
\providecommand{\decosix}{}
\providecommand{\starredbullet}{\ensuremath{\bullet}}
\pagestyle{empty}
\setlength{\parindent}{0pt}
\usepackage{adjustbox}\DeclareUnicodeCharacter{2236}{:}\DeclareUnicodeCharacter{2190}{\ensuremath{\leftarrow}}
\providecommand{\R}{\mathbb{R}}
\providecommand{\N}{\mathbb{N}}
\providecommand{\Z}{\mathbb{Z}}
\providecommand{\Q}{\mathbb{Q}}
\providecommand{\D}{\mathbb{D}}
\providecommand{\C}{\mathbb{C}}
\providecommand{\e}{{\mathrm{e}}}
\providecommand{\vect}[1]{\overrightarrow{#1}}
\providecommand{\vectt}[1]{\overrightarrow{\mathrm{#1}}}
\providecommand{\vvt}[1]{\overrightarrow{\mathrm{#1}}}
\providecommand{\np}[2][]{{#2~#1}}
\providecommand{\vv}[1]{{\overrightarrow{#1}}}
\renewcommand{\d}{{\mathrm{d}}}
\providecommand{\degres}{{^\circ}}
\providecommand{\degre}{{^\circ}}
\providecommand{\blue}{\color{blue}}
\providecommand{\red}{\color{red}}
\providecommand{\enskip}{\quad}
\providecommand{\llbracket}{\lbrack\!\lbrack}
\providecommand{\rrbracket}{\rbrack\!\rbrack}
\providecommand{\euro}{{\text{€}}}
\providecommand{\no}{{\text{n°}}}
\providecommand{\interval}[2]{{[#1;#2]}}
\providecommand{\intervFO}[2]{{[#1;#2[}}
\providecommand{\dotfill}{\ldots}

\providecommand{\ds}{\displaystyle}
\providecommand{\textding}[1]{#1}
\providecommand{\Oij}{\(\mathrm{O};\vec{\imath},\vec{\jmath}\)}
\providecommand{\Oijk}{\(\mathrm{O};\vec{\imath},\vec{\jmath},\vec{k}\)}
\providecommand{\Ouv}{\(\mathrm{O};\vec{u},\vec{v}\)}
\providecommand{\intervFF}[2]{\ensuremath{[#1;#2]}}
\renewcommand{\euro}{\ifmmode\text{\texteuro}\else\texteuro\fi}
\DeclareUnicodeCharacter{00B0}{\ifmmode^{\circ}\else\textdegree\fi}
\DeclareUnicodeCharacter{20AC}{\ifmmode\text{\texteuro}\else\texteuro\fi}
% Énoncés : Éducation nationale. Transcription et corrigés : APMEP, Denis Vergès et contributeurs. Découpage et classement : CoopMaths, Sébastien Lozano, Éric Elter et contributeurs. Adaptation de la présentation : Maths Explorer.
% CC BY-SA : https://coopmaths.fr/www/about/
% https://coopmaths.fr/alea/static/eam/2025/tex/eam_specifique_2025_06_sujet0v2_3.tex
\begin{document}
\noindent\textit{APMEP / CoopMaths. CC BY-SA. Présentation adaptée par Maths Explorer.}\par\medskip
\begin{center}
Les deux parties de cet exercice sont indépendantes.
\end{center}

\subsubsection*{Partie A}
Dans un lycée comptant 2\,000 élèves, on donne la répartition des effectifs suivant le sexe et le choix de la LV1.

\begin{center}
\begin{tabular}{|c|c|c|}
\hline
 & Fille & Garçon \\
\hline
Anglais & 712 & 728 \\
\hline
Autre LV1 & 288 & 272 \\
\hline
\end{tabular}
\end{center}

\begin{enumerate}
\item Un élève affirme $\textit{\og Dans ce lycée, il y a autant de filles que de garçons \fg}$. A-t-il raison ? Justifier.
\end{enumerate}

On choisit au hasard, de manière équiprobable, un élève dans ce lycée.

On considère les événements suivants :

$F~:~\textit{\og l'élève est une fille \fg}~;$ \\
$A~:~\textit{\og l'élève a choisi Anglais pour LV1 \fg}~.$

Dans les questions qui suivent, on donnera les résultats sous forme d'une fraction qu'il n'est pas demandé de simplifier.

\begin{enumerate}[resume]
\item Déterminer la probabilité de l'événement $A \cap F$.

\item Déterminer la probabilité de l'événement $A$ sachant que l'événement $F$ est réalisé.

\item Les événements $A$ et $F$ sont-ils indépendants ? Justifier.

\item On sait que l'élève choisi est un garçon. \\
On considère l'affirmation suivante :

\noindent\begin{minipage}{\linewidth}
\itshape\og La probabilité qu'il ait choisi Anglais pour LV1 est plus de trois fois plus grande que la probabilité qu'il n'ait pas choisi Anglais pour LV1 \fg.
\end{minipage}

Cette affirmation est-elle vraie ? Justifier.
\end{enumerate}

\subsubsection*{Partie B}
On dispose d'une pièce de monnaie truquée pour laquelle la probabilité d'obtenir pile lors d'un lancer est égale à $\dfrac{1}{4}$.

\begin{enumerate}
\item Déterminer la probabilité d'obtenir face.

\item On lance trois fois de suite cette pièce de monnaie, les trois lancers étant indépendants, et on note pour chaque lancer le résultat (pile ou face) obtenu.

\begin{enumerate}
\item Représenter la situation par un arbre de probabilités.

\item Quelle est la probabilité d'obtenir exactement une fois pile lors de ces trois lancers ?

\item Quelle est la probabilité de ne jamais obtenir pile ?
\end{enumerate}
\end{enumerate}

\end{document}
